
\section{Ontology}
\label{sect:ontology}

To meet \ref{sect:minimal-requirements} we start by assuming an observer exist. This is the first principle from which
the entire theory is derived. 

We do not assume finiteness.

If reality were fundamentally finite, one could immediately ask: finite by how much? Why n bits rather than n+1?
What principle fixed that particular bound?
Any fixed finite limit appears contingent and therefore calls for a deeper explanatory layer.
For this reason, we do not posit a fundamentally finite substrate.

We do not assume information.

Information presupposes distinguishability. At its simplest, this reduces to binary distinction: a bit, a difference between alternatives.
But declaring reality to be fundamentally informational or digital is itself a substantive ontological commitment.
A discrete ontology is conceivable, but so is a continuous one. Nothing prior licenses choosing one baseline over the other.
Therefore, we do not presuppose that reality is fundamentally finite, discrete, informational, mathematical, or governed by prior, external laws.

We denote this unrestricted, pre-interpretive totality ($\mathbf{U}$). We do not claim to know what it is.

Only that no specific formatting principle is granted ontological privilege a priori.



